\begin{textsample}{NEG Dim 7 – global\_south\_2025\_02 – Score -1.00 – global\_south\_2025\_02/121543314.txt}  \label{ex:f7_neg_039}
Gaza/Jerusalem : Palestinian armed groups on Thursday handed over the bodies of four Israeli hostages to the International Committee of the Red Cross ( ICRC ) team in Khan Younis in southern Gaza Strip, said Palestinian sources.
Marking the first repatriation of deceased hostages under an ongoing ceasefire agreement, the handover occurs near \textbf{Bani} Suhaila cemetery during a military parade attended by members of al-Qassam Brigades, the armed wing of Hamas and other Palestinian factions.
It featured a large crowd of masked gunmen positioned around a platform where the coffins were placed.
After the ICRC team arrived, a militant of al-Qassam Brigades lifted a black curtain, revealing four black wooden coffins adorned with pictures of the deceased.
Hamas said in a press statement that the bodies belonged to Shiri Bibas, her two children, Kfir and Ariel Bibas, and an elderly man Oded Lifshitz, who were killed in the Gaza Strip during the conflict with Israel.
The process formally began with a Hamas representative signing a transfer document alongside an ICRC official. @ @ @ @ @ @ @ @ @ @ in Gaza, will be transferred to the National Centre of Forensic Medicine in Tel Aviv for formal identification, Israel ' s Prime Minister ' s Office said in a statement."At the conclusion of the identification process, an official notification will be given to the families,"the office said.
The four were taken hostage during a Hamas attack on Israel on October 7, 2023.
The Mujahideen Brigades, the armed wing of the Palestinian Mujahideen Movement, claimed responsibility for capturing the Bibas family during the attack, while the Al-Quds Brigades, the armed wing of the Palestinian Islamic Jihad movement, said that they had detained Lifshitz.
Both factions asserted that all Israeli captives were alive until Israeli airstrikes targeted their detention sites.
Following the handover, Hamas said in a press statement that the return of the bodies was carried out for humanitarian reasons despite what it described as"ongoing violations"by the Israeli military against Palestinian prisoners, Xinhua news agency reported.
Accusing Israel of committing a"new crime @ @ @ @ @ @ @ @ @ @ full responsibility for their deaths, as the captives were allegedly killed in direct Israeli airstrikes alongside the Palestinian captors.
The armed group also accused the Israeli government of repeatedly obstructing hostage-for-prisoner exchange negotiations, prolonging the captives ' detention and ultimately leading to their deaths in military actions. ( IANS )

% matched lemmas: bani
\end{textsample}
